\scanpage{197}

\deleted{\unsure{Abordons|Considérons} en particulier le cas où $X\to S$ est propre. \unsure{Donc on a|Dans ce cas} $\underline{\unsure{W|\omega}}=f_*(\Omega^1_{X/S})$ faisceau coh. sur $S$, et quitte à restreindre $S$, on peut supposer que}

\[
\deleted{\underline{\underline{W}}=f_*(Z^1_{X/S})=\operatorname{Ker}\bigl(f_*(\Omega_{X/S})\longrightarrow f_*(\Omega^2_{X/S})\bigr).}
\]

Dans le cas où $X\to S$ est un schéma abélien, introduisant le schéma abélien dual $A$, et son extension vectorielle universelle $E$, la question \deleted{\unsure{se|est}} équivaut à la suivante : \deleted{Soit $\varphi$ une section}

\margin{cette traduction vaut, plus gén., si $X\to S$ propre et lisse}

% vertical rule beside c)

c) Soit $\deleted{\varphi}\omega$ une section de $E$ sur $S$. Supposons que pour tout pt \uline{fermé} $s\in S$, $\omega(s)\in A_s^{\unsure{(p/k(s))|(p/k_s)|(p(k(s)))}}$ \inserted{$(=pE_s)$}, où $p=\mathrm{car}\,k(s)$. Alors la \deleted{\unsure{l'image de|image de}}\inserted{section} $\omega$ de $\Gamma(S,A_s)$ est-elle d'\uline{ordre fini} ? [NB. a) signifie que l'image de $\omega$ dans $\Gamma(S,A_S)$ est d'ordre fini, b) signifie que si on \inserted{suppose de plus que} $\omega$ est une section de $V=\mathrm{Ker}(E\to A)$, alors $\omega$ est nulle].

Ceci suggère la question

d) Soit $S$ de type fini sur $\mathbb Z$, intègre, \inserted{de pt gén $\eta$}, dominant $\mathrm{Spec}\,\mathbb Z$, \deleted{\unsure{Soit}} $E$ un schéma en groupes \deleted{\unsure{de loc}} de t.f. sur $S$ (lisse connexe si on veut\unsure{ !|}), $\omega$ une section de $E$. Supposons que pour tt pt fermé $s\in S$, $\omega(s)$ appartienne \deleted{au ss-gpe} à un sous-gpe de $E_s$ qui est \uline{propre} \inserted{\uline{connexe}} sur $k(s)$ [resp. \inserted{connexe} de rang unipotent nul, resp. \inserted{\uline{connexe}} affine, resp. \inserted{connexe} unipotent] alors $\omega(\eta)$ appartient à un ss-gpe de $E_\eta$ qui est \uline{propre}

% archive sheet number: 196
